\scanpage{19}
\[
(8)\quad\left\{\begin{aligned}
\varepsilon_0&=l_0^{\deleted{\unsure{-1}}}\sigma_0l_1^{\deleted{\unsure{-1}}}\\
\varepsilon_1&=l_1^{\deleted{\unsure{-1}}}\sigma_1l_\infty^{\deleted{\unsure{-1}}}\\
\varepsilon_\infty&=l_\infty\sigma_\infty l_0^{\deleted{\unsure{-1}}}
\end{aligned}\right.
\]
\margin{\uline{NB}\quad $\rho\varepsilon_i\rho^{-1}=\varepsilon_{\rho i}$ i.e.
\[
\left\{\begin{aligned}
\rho\varepsilon_0\rho^{-1}&=\varepsilon_1\\
\rho\varepsilon_1\rho^{-1}&=\varepsilon_\infty\\
\rho\varepsilon_\infty\rho^{-1}&=\varepsilon_0
\end{aligned}\right.
\]}
On a bien
\[
\varepsilon_0^2=l_0\underbrace{\sigma_0l_1l_0\sigma_0}_{\substack{\sigma_0(l_1l_0)=\sigma_0(l_1)\sigma_0(l_0)\\=l_\infty(l_1l_\infty)^{-1}=l_1^{-1}}}l_{\unsure{1}}=l_0l_\infty l_\infty^{-1}l_1^{-1}l_1=l_0\qquad\text{OK}
\]
(\uline{NB} Les $\varepsilon_i$ ici sont les $\varepsilon_i^{-1}$ de tantôt ! \deleted{\unsure{il faudrait}} \inserted{qu'il faudra} \deleted{\unsure{ici}} donc noter $\varepsilon'_i\ldots$)

Ici on voit donc qu'on a un syst. de gén. et relations \inserted{de $\mathcal E^+$} $l_0,l_1,l_\infty,\varepsilon_0,\varepsilon_1,\varepsilon_\infty,\rho$ \deleted{de $\mathcal E^+$}, les relations
\[
\left\{\begin{gathered}
l_\infty l_1l_0=1\\
\varepsilon_0^2=l_0,\quad\varepsilon_1^2=l_1,\quad\varepsilon_\infty^2=l_\infty\\
\varepsilon_0\varepsilon_1=\rho\deleted{\unsure{l_\infty}}(\ldots),\quad\varepsilon_1\varepsilon_\infty=\ldots,\quad\varepsilon_\infty\varepsilon_0=\ldots,\quad\rho^3=1
\end{gathered}\right.
\]
où il faut compléter l'expression de $\varepsilon_0\varepsilon_1$ etc \inserted{comme} produits de $\rho$ par des expressions en les $l_i$.
\deleted{\[
\varepsilon_0\varepsilon_1=l_0\sigma_1l_1
\]}
\[
\varepsilon_0\varepsilon_1=l_0\sigma_0l_1l_1\sigma_1l_\infty
=l_0\underbrace{\sigma_0l_1^2\sigma_0^{-1}}_{l_\infty^2}\underbrace{\sigma_0\sigma_1}_{\rho l_0}l_\infty
=l_0l_\infty^2\rho\unsure{l_0}l_\infty
\]
donc finalement
\[
(9)\quad\left\{\begin{aligned}
\varepsilon_0\varepsilon_1\deleted{\unsure{\varepsilon}}&=l_0l_\infty^2\rho\unsure{l_0}l_\infty\\
\varepsilon_1\varepsilon_\infty&=l_1l_0^2\rho\unsure{l_1}l_0\\
\varepsilon_\infty\varepsilon_0&=l_\infty l_1^2\rho\unsure{l_\infty}l_1
\end{aligned}\right.
\]
Eliminant les $l_i$ comme $\varepsilon_i^2$, on trouve des gén. $\varepsilon_0,\varepsilon_1,\varepsilon_\infty,\rho$, avec relations de \unsure{définition}
\mspage{18}
